% ch12-decision — from model output to investment decisions

\section{Closing the loop with Chapter 1}\label{ch12:sec:loop}

\Cref{ch:question} defined worth by the value-of-information identity
\cref{eq:ch01:voi}: a model output matters only through the action it
improves. The intervening chapters built the outputs --- a sector-level
ground intensity (\cref{ch:poisson,ch:hawkes,ch:censoring}), a
firm-level conditional ranker with an outside option
(\cref{ch:ranking,ch:universe}), horizon event probabilities
(\cref{ch:buckets}), all under the evaluation constitution of
\cref{ch:gof} and the covariate honesty of \cref{ch:stale}. This chapter
specifies the \emph{decision layer}: the small set of composed,
consumable deliverables, the decision rules that consume them, and the
feedback loop that measures whether the whole apparatus is earning its
keep. The layer is deliberately thin --- it contains no new statistical
estimation, only composition, thresholds, and reporting --- because every
quantity it consumes arrives already calibrated and already
uncertainty-qualified; anything thicker would smuggle modeling into a
place with no likelihood to check it.

\section{The three deliverables}\label{ch12:sec:deliverables}

\subsection{The sourcing list (T1 + T3)}

The headline deliverable is a per-sector ranked list, refreshed weekly:
for each sector $s$ at time $\tau$, the top-$K$ tracked firms by
conditional next-deal probability $\hat p(i \mid s, \tau)$
(\cref{ch:ranking}), each annotated with its horizon probabilities
$\hat p_i(\tau, \horizon)$ for $\horizon \in \{$6M, 1Y, 2Y$\}$
(\cref{ch:buckets}), its staleness age $\staleage_i(\tau)$ and last-known
covariates (honesty about the information state, \cref{ch:stale}), and
the sector's calibrated newcomer share $\hat p(\lozenge_s)$
(\cref{ch:universe}). The two probability columns answer different
questions and must not be merged: the ranker column is \emph{relative}
(who, given a deal), the horizon column is \emph{absolute} (whether,
within $\horizon$). A firm can top the ranker while all horizon
probabilities are low because the sector itself is quiet; the pairing is
exactly what tells the desk to watch rather than act.

\begin{proposition}[Optimality of the top-$K$ sourcing rule]\label{prop:ch12:topk}
Suppose diligence capacity allows opening at most $K$ files in sector
$s$, each opened file on firm $i$ yields expected utility $b \cdot
\ind{i \text{ raises within } \horizon}$ minus cost $c$, with $b, c$
constant across firms, and outcomes are as modeled. Then the expected-utility-maximizing
selection is the top-$K$ firms by $\hat p_i(\tau,\horizon)$, and opening
the $k$-th file is worthwhile iff $\hat p_{(k)}(\tau,\horizon) \ge c/b$.
\end{proposition}

\begin{proof}
Expected utility of a selection $A$, $|A| \le K$, is $\sum_{i \in A}
\big(b\, \hat p_i(\tau,\horizon) - c\big)$, additive in members; it is
maximized by taking the firms with the largest positive terms, which are
the top-ranked ones by $\hat p_i$, truncated at $K$ or at the threshold
$\hat p_i \ge c/b$, whichever binds first. The threshold form is the
classical cost--benefit decision boundary for calibrated probabilities;
the argument fails exactly when $\hat p$ is uncalibrated, which is why
\cref{ch:buckets} makes calibration mandatory, and when utilities
interact across firms (portfolio effects), which is out of scope by
\cref{ch12:sec:scope}.
\end{proof}

The ranking-only analogue --- expected recall of the next deals is
maximized by ranking on conditional probability --- was proved as the
top-$K$ optimality lemma of \cref{ch:ranking}. Together the two results
say: \emph{under the model, the list is the decision}; everything else
is threshold arithmetic supplied by the desk's own $c/b$.

\subsection{Sector pacing (T2)}

For capital-deployment pacing the deliverable is the predictive
distribution of sector deal counts over the next 4--13 weeks: mean paths
with predictive intervals, produced by simulating the fitted ground
process forward (\cref{ch:hawkes,ch:censoring} simulators; scenario or
bootstrapped macro paths per \cref{ch:data}), with dispersion honestly
propagated (negative-binomial noise where adopted, \cref{ch:poisson}).
Every pacing chart carries two inherited caveats as fixed footnotes: the
fitted branching mass is an upper bound on contagion
(\cref{ch:hawkes}), and timing resolution is bounded below by the honest
resolution of \cref{ch:censoring} --- no intra-week timing claims.

\subsection{Deal-terms context (T4)}

Mark regressions --- round size, up/down/flat probability, valuation
step conditional on a deal --- contextualize a candidate: the same
gradient-boosting stack as \cref{ch:buckets} applied to deal-level
targets from \cref{app:columns}, with staleness features and the
undisclosed-size selection handled per the data dictionary's imputation
rules (\code{deal\_size\_status} is itself informative and enters as a
feature, never silently imputed away). Terms models are context, not
gates: they refine the utility $b$ in \cref{prop:ch12:topk}, firm by
firm, when the desk wants probability-weighted expected terms rather
than a flat $b$.

\section{Scope: what the decision layer refuses to do}\label{ch12:sec:scope}

\begin{decision}{D12.1 --- no silent portfolio optimization}
The decision layer outputs ranked lists, calibrated probabilities, and
predictive count distributions. It does \emph{not} output buy/pass
verdicts, position sizes, or fund-level allocations. Rationale:
portfolio-level utility involves correlations, capacity, and mandate
constraints that live outside the data this document models; any such
optimization must be a separate, owned system consuming our calibrated
outputs. Reversal trigger: a stated fund-level objective with its own
evaluation protocol.
\end{decision}

\begin{decision}{D12.2 --- probabilities travel with their caveats}
Every number leaving the system carries: the model version, the
evaluation-era metrics of \cref{ch:gof}, the risk-set definition
(\cref{ch:ranking}), the staleness summary of the inputs
(\cref{ch:stale}), and the newcomer share (\cref{ch:universe}). The
reporting discipline of \cref{ch:gof} --- no hazard without its risk-set
rule, no branching ratio without its upper-bound caveat, no newcomer
probability sold as an entrant probability --- binds the decision layer
most of all, because it is the layer other people quote. Reversal
trigger: none; this is charter.
\end{decision}

\section{Measuring the value of information}\label{ch12:sec:voi}

Offline metrics are proxies. The loop closes only by measuring
\cref{eq:ch01:voi} on the workflow itself.

\begin{protocol}[VoI measurement]\label{prot:ch12:voi}
(a) \emph{Shadow deployment}: for at least one full quarter, sourcing
lists are produced and frozen weekly but not shown to the desk; realized
deals score them retrospectively (recall@$K$, lift over the desk's
contemporaneous shortlists). (b) \emph{Assisted period}: lists shown;
the desk logs which opened files originated from the list vs.\ other
channels. (c) The VoI estimate is the difference in hit rates per opened
file between assisted and baseline periods, era-blocked as in
\cref{ch:gof}; secular market shifts are partially controlled by the
shadow-period comparison. (d) Attribution is reported with its
confounds stated --- this is an observational estimate, not a randomized
one; if the desk's capacity permits, randomize \emph{which sectors} get
assisted lists per quarter for a cleaner contrast.
\end{protocol}

\section{Monitoring, drift, and retraining cadence}\label{ch12:sec:monitoring}

Models decay: coverage drifts (\cref{ch:universe}), macro regimes shift,
and the 2021--2023 funding cycle sits inside our evaluation eras as a
standing warning. The cadence is quarterly re-evaluation against the
rolling-origin protocol of \cref{ch:gof}, with three tripwires, each
mapped to its owning chapter: (i) calibration drift on horizon
probabilities beyond the tolerance band of the reliability protocol
(\cref{ch:buckets}) triggers recalibration before retraining; (ii) a
trend in \code{off\_universe\_share} or in the Good--Turing missing mass
(\cref{ch:universe}) triggers the coverage audit before any modeling
response; (iii) attenuation-gap growth in covariate effects
(\cref{ch:stale}'s gap test) escalates the staleness remedy ladder.
Retraining without a fired tripwire is discouraged: it churns the
desk's mental model of the lists for no measured gain.

\begin{codehooks}
The decision layer is a reporting module, not a modeling module:
\modrepo{decision/} (to be written) composes existing outputs ---
\code{ranker} scores from \modrepo{sector\_ranker.py} (extended with the
outside option), horizon probabilities from \modrepo{models/xgb\_horizon.py}
(planned, \cref{ch:buckets}), simulated count fans from the ground-process
simulators, and terms context --- into versioned weekly artifacts
(per-sector JSON + rendered sheets). The VoI protocol needs a tiny
\modrepo{decision/log.py} for desk-interaction logging. Tripwire metrics
reuse the \modrepo{eval/} harness of \cref{ch:gof}; nothing in this layer
may compute a probability itself.
\end{codehooks}
